\section{Related Work}
\label{sec:related}

\subsection{Few-Shot Learning}
\label{sec:related_fewshot}

\paragraph{Meta-learning}
Few-shot learning asks a model to solve a new task from a few examples, and meta-learning makes this ability the training objective~\cite{hospedales2021survey}.
Metric-based methods compare queries to a support set in a learned embedding space~\cite{vinyals2016matching,snell2017proto}.
Optimization-based methods instead learn an initialization from which a few gradient steps suffice.
MAML~\cite{finn2017maml} differentiates through the inner loop, its first-order variant drops the second-order terms, and Reptile~\cite{nichol2018reptile} replaces the meta-gradient by an interpolation toward the weights reached after $k$ inner steps, which decouples the inner horizon from the memory footprint.
ANIL~\cite{raghu2020anil} shows that much of MAML's benefit on few-shot classification comes from feature reuse with a rapidly adapting head, and \cite{antoniou2019howto} documents the instability of MAML-style training.
We use Reptile, but with an inner loop that is the full post-training recipe of a 1.4B-parameter trainable set rather than a few steps on a small network, and we rely on the fact that its update is a convex combination of weights, which lets the standard sharded training stack run the inner loop.

\paragraph{Few-shot imitation and data-efficient fine-tuning}
In robotics, one-shot imitation learning conditions a policy on a demonstration~\cite{duan2017oneshot,jang2022bcz}, and meta-imitation learning applies MAML to visuomotor policies so that a single demonstration adapts them~\cite{finn2017mil,yu2018daml}; skill-based methods adapt a pre-trained skill prior to a few demonstrations~\cite{hakhamaneshi2022fist}.
These works meta-train small policies on tens of tasks of one robot.
A complementary line improves data efficiency through pre-trained visual representations~\cite{nair2022r3m,majumdar2023vc1} or by fine-tuning generalist policies~\cite{octo2024,kim2024openvla,black2024pi0,intelligence2025pi05,shukor2025smolvla}, typically with parameter-efficient adapters~\cite{hu2022lora,kim2025oft}.
Their fine-tuning sets still contain tens to hundreds of demonstrations per task; the regime of three demonstrations per task with a billion-parameter model has, to our knowledge, not been studied.
Our method is orthogonal to the choice of adapter and recipe: it changes the point from which post-training starts, not how post-training proceeds.

\subsection{World Action Models}
\label{sec:related_wam}

World models learn to predict the consequences of actions~\cite{ha2018worldmodels,hafner2020dreamer}.
With video generation this idea has scaled to pixels.
UniPi~\cite{du2023unipi} plans by generating a video and recovers actions with an inverse dynamics model; GR-1 and GR-2~\cite{wu2024gr1,cheang2024gr2} pre-train a video generator and then predict actions alongside future frames; VPP~\cite{hu2025vpp} feeds the representations of a video diffusion model to a policy; Genie~\cite{bruce2024genie} and LAPA~\cite{ye2025lapa} learn latent actions from unlabeled videos; DreamGen~\cite{jang2025dreamgen} uses a post-trained video world model to synthesize training trajectories; and V-JEPA~2~\cite{assran2025vjepa2} plans in a learned latent space.
A growing family couples video and action generation in one network, which we refer to as world action models: UWM~\cite{zhu2025uwm} couples video and action diffusion, WorldVLA~\cite{cen2025worldvla} autoregresses over both modalities, Genie Envisioner~\cite{liao2025genieenvisioner} builds a policy on a video world model, and Cosmos~\cite{cosmos1,cosmos3} unifies language, video, and action generation in a mixture-of-transformers backbone~\cite{liang2024mot,zhou2024transfusion}.
Vision-language-action models~\cite{brohan2023rt1,zitkovich2023rt2,kim2024openvla,black2024pi0,octo2024,nvidia2025groot} are the action-only relatives of these models, trained on cross-embodiment collections~\cite{oneill2024oxe,khazatsky2024droid}.
All of these models are post-trained on target data with tens to hundreds of demonstrations per task; we study how to make this post-training work with three.
We use \base~\cite{cosmos3} as our WAM, but the meta pre-training stage only assumes a post-training recipe and a set of source robot datasets, and is agnostic to the backbone.
